\documentclass[11pt]{beamer}

\usetheme{Madrid}
\usecolortheme{default}
\definecolor{unicaucaverde}{RGB}{0,90,60}
\setbeamercolor{structure}{fg=unicaucaverde}
\setbeamertemplate{navigation symbols}{}

\usepackage[utf8]{inputenc}
\usepackage[spanish]{babel}
\usepackage{amsmath,amssymb}

\newcommand{\R}{\mathbb{R}}
\newcommand{\vv}[1]{\mathbf{#1}}

\title[Semana 4: Matriz Inversa y Transpuesta]{Matrices, matriz inversa y transpuesta}
\subtitle{Álgebra Lineal para Ingeniería --- Semana 4 de 16}
\author{Universidad del Cauca}
\date{2026}

\begin{document}

\begin{frame}
  \titlepage
\end{frame}

\begin{frame}{Semana 4: dónde estamos}
  \begin{itemize}
    \item Unidad 2 de 8: Vectores y matrices (Grossman, cap. 2).
    \item Subtemas 2.3, 2.4 y 2.5 --- RA2.3, RA2.4, RA2.5.
    \item 4 horas: Forma matricial $A\vv{x}=\vv{b}$ $\to$ Matriz inversa $A^{-1}$ $\to$ Resolución de sistemas $\to$ Matriz transpuesta $A^T$.
  \end{itemize}
  \tableofcontents
\end{frame}

\section{Matrices y sistemas de ecuaciones lineales}

\begin{frame}{Forma matricial de un sistema}
  \begin{block}{Sistema de $m$ ecuaciones con $n$ incógnitas}
    \[
      A\vv{x} = \vv{b}
    \]
    \begin{itemize}
      \item $A \in \R^{m\times n}$: matriz de coeficientes.
      \item $\vv{x} \in \R^n$: vector de incógnitas.
      \item $\vv{b} \in \R^m$: vector de términos independientes.
    \end{itemize}
  \end{block}
  \vfill
  \begin{block}{Combinación lineal de columnas}
    $A\vv{x} = x_1 \vv{a}_1 + x_2 \vv{a}_2 + \cdots + x_n \vv{a}_n = \vv{b}$.\\
    \textbf{Consistencia:} El sistema es consistente si y solo si $\vv{b}$ es combinación lineal de las columnas de $A$.
  \end{block}
\end{frame}

\section{Inversa de una matriz cuadrada}

\begin{frame}{Definición de matriz inversa}
  \begin{block}{Matriz invertible (no singular)}
    Sea $A$ de orden $n\times n$. Si existe una matriz $B$ tal que
    \[
      A B = B A = I_n,
    \]
    entonces $B = A^{-1}$ es la \textbf{inversa única} de $A$.
  \end{block}
  \vfill
  Si no existe tal matriz, $A$ se dice \textbf{singular}.
\end{frame}

\begin{frame}{Inversa de matriz $2\times 2$}
  \begin{block}{Fórmula explícita}
    Para $A = \begin{pmatrix} a & b \\ c & d \end{pmatrix}$, si el determinante $\Delta = ad - bc \neq 0$:
    \[
      A^{-1} = \frac{1}{ad - bc} \begin{pmatrix} d & -b \\ -c & a \end{pmatrix}.
    \]
  \end{block}
  \vfill
  \textbf{Ejemplo:} Para $A = \begin{pmatrix} 2 & 5 \\ 1 & 3 \end{pmatrix}$:
  $\Delta = 6 - 5 = 1 \implies A^{-1} = \begin{pmatrix} 3 & -5 \\ -1 & 2 \end{pmatrix}$.
\end{frame}

\begin{frame}{La Matriz Identidad $I_n$}
  \begin{block}{Definición y notación}
    $I_n = (\delta_{ij})_{n\times n}$ con $\delta_{ij}=1$ si $i=j$ y $\delta_{ij}=0$ si $i\neq j$.
    \[
      I_2 = \begin{pmatrix} 1&0\\0&1 \end{pmatrix}, \qquad I_3 = \begin{pmatrix} 1&0&0\\0&1&0\\0&0&1 \end{pmatrix}.
    \]
  \end{block}
  \vfill
  \begin{block}{Propiedades clave}
    \begin{itemize}
      \item \textbf{Elemento neutro:} $A I_n = A$ e $I_m A = A$.
      \item \textbf{Transformación idéntica:} $I_n \vv{x} = \vv{x}$ para todo $\vv{x}\in\R^n$.
      \item \textbf{Auto-inversa:} $I_n^{-1} = I_n$.
    \end{itemize}
  \end{block}
\end{frame}

\begin{frame}{Fundamentación del Método de Gauss-Jordan}
  \begin{block}{¿Por qué funciona $[A \mid I_n] \to [I_n \mid A^{-1}]$?}
    Cada operación elemental de fila equivale a multiplicar a la izquierda por una \textbf{matriz elemental} $E_k$.
    \vfill
    Si $A$ es no singular, existe una secuencia de operaciones tales que:
    \[
      (E_k \cdots E_2 E_1) A = I_n \implies E_k \cdots E_2 E_1 = A^{-1}.
    \]
    \vfill
    Aplicando la misma secuencia sobre $I_n$:
    \[
      (E_k \cdots E_2 E_1) I_n = A^{-1} I_n = A^{-1}.
    \]
  \end{block}
\end{frame}

\begin{frame}{Algoritmo de Gauss-Jordan paso a paso}
  \begin{enumerate}
    \item \textbf{Aumentar:} Construir $[A \mid I_n]$ de dimensión $n\times 2n$.
    \item \textbf{Pivoteo ($j=1\dots n$):}
      \begin{itemize}
        \item Si $a_{jj}=0$, intercambiar filas $R_j \leftrightarrow R_k$.
        \item Normalizar pivote a $1$: $R_j \leftarrow \frac{1}{a_{jj}} R_j$.
        \item Hacer ceros en columna $j$ (arriba y abajo): $R_r \leftarrow R_r - a_{rj} R_j$.
      \end{itemize}
    \item \textbf{Resultado:} Si la mitad izquierda llega a $I_n$, la derecha es $A^{-1}$:
      \[
        [A \mid I_n] \xrightarrow{\text{Gauss-Jordan}} [I_n \mid A^{-1}].
      \]
  \end{enumerate}
\end{frame}

\begin{frame}{Propiedades de la matriz inversa}
  \begin{itemize}
    \item $(A^{-1})^{-1} = A$.
    \item $(AB)^{-1} = B^{-1} A^{-1}$ \quad (\textbf{¡orden invertido!}).
    \item $(cA)^{-1} = \frac{1}{c} A^{-1}$ \quad ($c \neq 0$).
  \end{itemize}
  \vfill
  \begin{block}{Solución de sistemas con $A^{-1}$}
    Si $A$ es invertible, la única solución de $A\vv{x} = \vv{b}$ es:
    \[
      \vv{x} = A^{-1}\vv{b}.
    \]
  \end{block}
\end{frame}

\section{Transpuesta de una matriz}

\begin{frame}{Definición y propiedades de $A^T$}
  \begin{block}{Matriz transpuesta}
    Dada $A=(a_{ij})_{m\times n}$, su transpuesta $A^T=(a_{ji})_{n\times m}$ intercambia filas por columnas.
  \end{block}
  \vfill
  \begin{block}{Propiedades}
    \begin{itemize}
      \item $(A^T)^T = A$.
      \item $(A+B)^T = A^T + B^T$.
      \item $(cA)^T = c A^T$.
      \item $(AB)^T = B^T A^T$ \quad (\textbf{¡orden invertido!}).
      \item $(A^T)^{-1} = (A^{-1})^T$.
    \end{itemize}
  \end{block}
\end{frame}

\begin{frame}{Matrices simétricas y antisimétricas}
  Para $A$ cuadrada $n\times n$:
  \begin{itemize}
    \item \textbf{Simétrica:} $A^T = A$ \quad ($a_{ij} = a_{ji}$).
    \item \textbf{Antisimétrica:} $A^T = -A$ \quad ($a_{ij} = -a_{ji}$, diagonal en cero $a_{ii}=0$).
  \end{itemize}
  \vfill
  \begin{block}{Descomposición única}
    Toda matriz cuadrada $A$ se puede escribir como:
    \[
      A = \underbrace{\frac{A + A^T}{2}}_{\text{Simétrica } S} + \underbrace{\frac{A - A^T}{2}}_{\text{Antisimétrica } K}.
    \]
  \end{block}
\end{frame}

\section{Ejercicios resueltos (síntesis)}

\begin{frame}{Ejercicio 1: Forma matricial}
  \textbf{Sistema:}
  $\begin{cases} x_1 + 2x_2 - x_3 = 5 \\ 3x_1 - x_2 + 2x_3 = 3 \\ -x_1 + x_2 + x_3 = -2 \end{cases}$
  \vfill
  \textbf{Forma matricial:}
  \[
    \begin{pmatrix} 1 & 2 & -1 \\ 3 & -1 & 2 \\ -1 & 1 & 1 \end{pmatrix} \begin{pmatrix} x_1 \\ x_2 \\ x_3 \end{pmatrix} = \begin{pmatrix} 5 \\ 3 \\ -2 \end{pmatrix}.
  \]
  Verificación con $\vv{x}=(2,1,-1)^T$:
  $1(2)+2(1)-(-1)=5$, $3(2)-1+2(-1)=3$, $-2+1-1=-2$. \checkmark
\end{frame}

\begin{frame}{Ejercicio 2: Inversa $2\times 2$ y solución de sistema}
  Dada $A = \begin{pmatrix} 4 & 2 \\ 3 & 2 \end{pmatrix}$, hallar $A^{-1}$ y resolver $A\vv{x} = \begin{pmatrix} 10 \\ 8 \end{pmatrix}$.
  \vfill
  \begin{enumerate}
    \item $\det(A) = 4(2)-2(3) = 2 \neq 0$.
    \item $A^{-1} = \frac{1}{2}\begin{pmatrix} 2 & -2 \\ -3 & 4 \end{pmatrix} = \begin{pmatrix} 1 & -1 \\ -3/2 & 2 \end{pmatrix}$.
    \item $\vv{x} = A^{-1}\vv{b} = \begin{pmatrix} 1 & -1 \\ -3/2 & 2 \end{pmatrix}\begin{pmatrix} 10 \\ 8 \end{pmatrix} = \begin{pmatrix} 2 \\ 1 \end{pmatrix}$.
  \end{enumerate}
\end{frame}

\begin{frame}{Ejercicio 3: Inversa $3\times 3$ por Gauss-Jordan}
  Hallar $A^{-1}$ para $A = \begin{pmatrix} 1 & 0 & 2 \\ 2 & -1 & 3 \\ 4 & 1 & 8 \end{pmatrix}$:
  \[
    \left(\begin{array}{ccc|ccc} 1&0&2&1&0&0 \\ 2&-1&3&0&1&0 \\ 4&1&8&0&0&1 \end{array}\right)
    \xrightarrow{\text{Gauss-Jordan}}
    \left(\begin{array}{ccc|ccc} 1&0&0&-11&2&2 \\ 0&1&0&-4&0&1 \\ 0&0&1&6&-1&-1 \end{array}\right).
  \]
  \vfill
  \[
    A^{-1} = \begin{pmatrix} -11 & 2 & 2 \\ -4 & 0 & 1 \\ 6 & -1 & -1 \end{pmatrix}.
  \]
\end{frame}

\begin{frame}{Ejercicio 4: Propiedad $(AB)^{-1} = B^{-1} A^{-1}$}
  Para $A = \begin{pmatrix} 1 & 2 \\ 0 & 1 \end{pmatrix}$ y $B = \begin{pmatrix} 3 & 1 \\ 2 & 1 \end{pmatrix}$:
  \begin{itemize}
    \item $AB = \begin{pmatrix} 7 & 3 \\ 2 & 1 \end{pmatrix} \implies (AB)^{-1} = \begin{pmatrix} 1 & -3 \\ -2 & 7 \end{pmatrix}$.
    \item $A^{-1} = \begin{pmatrix} 1 & -2 \\ 0 & 1 \end{pmatrix}$, \quad $B^{-1} = \begin{pmatrix} 1 & -1 \\ -2 & 3 \end{pmatrix}$.
    \item $B^{-1} A^{-1} = \begin{pmatrix} 1 & -1 \\ -2 & 3 \end{pmatrix}\begin{pmatrix} 1 & -2 \\ 0 & 1 \end{pmatrix} = \begin{pmatrix} 1 & -3 \\ -2 & 7 \end{pmatrix}$.
  \end{itemize}
  \vfill
  ¡Verificado! Coinciden exactamente.
\end{frame}

\begin{frame}{Ejercicio 5: Transposición y matriz simétrica $AA^T$}
  Sea $A = \begin{pmatrix} 1 & -1 & 3 \\ 2 & 0 & 4 \end{pmatrix}$.
  \begin{itemize}
    \item $A^T = \begin{pmatrix} 1 & 2 \\ -1 & 0 \\ 3 & 4 \end{pmatrix}$.
    \item $M = AA^T = \begin{pmatrix} 1 & -1 & 3 \\ 2 & 0 & 4 \end{pmatrix} \begin{pmatrix} 1 & 2 \\ -1 & 0 \\ 3 & 4 \end{pmatrix} = \begin{pmatrix} 11 & 14 \\ 14 & 20 \end{pmatrix}$.
    \item $M^T = \begin{pmatrix} 11 & 14 \\ 14 & 20 \end{pmatrix} = M$ (es simétrica).
  \end{itemize}
\end{frame}

\begin{frame}{Ejercicio 6: Descomposición $A = S + K$}
  Descomponer $A = \begin{pmatrix} 2 & 5 \\ 1 & 4 \end{pmatrix}$ en simétrica $S$ y antisimétrica $K$:
  \begin{itemize}
    \item $S = \frac{1}{2}(A + A^T) = \frac{1}{2}\begin{pmatrix} 4 & 6 \\ 6 & 8 \end{pmatrix} = \begin{pmatrix} 2 & 3 \\ 3 & 4 \end{pmatrix}$ (simétrica).
    \item $K = \frac{1}{2}(A - A^T) = \frac{1}{2}\begin{pmatrix} 0 & 4 \\ -4 & 0 \end{pmatrix} = \begin{pmatrix} 0 & 2 \\ -2 & 0 \end{pmatrix}$ (antisimétrica).
    \item $S + K = \begin{pmatrix} 2 & 3 \\ 3 & 4 \end{pmatrix} + \begin{pmatrix} 0 & 2 \\ -2 & 0 \end{pmatrix} = \begin{pmatrix} 2 & 5 \\ 1 & 4 \end{pmatrix} = A$. \checkmark
  \end{itemize}
\end{frame}

\begin{frame}{Ejercicio 7: Método de la Identidad $[A \mid I_3] \to [I_3 \mid A^{-1}]$}
  Dada $A = \begin{pmatrix} 1 & 2 & 1 \\ 2 & 5 & 4 \\ 1 & 1 & 0 \end{pmatrix}$, hallar $A^{-1}$ reduciendo $[A \mid I_3]$:
  \[
    \left(\begin{array}{ccc|ccc} 1&2&1&1&0&0 \\ 2&5&4&0&1&0 \\ 1&1&0&0&0&1 \end{array}\right)
    \xrightarrow{F_2-2F_1, F_3-F_1}
    \left(\begin{array}{ccc|ccc} 1&2&1&1&0&0 \\ 0&1&2&-2&1&0 \\ 0&-1&-1&-1&0&1 \end{array}\right)
  \]
  \[
    \xrightarrow{F_3+F_2}
    \left(\begin{array}{ccc|ccc} 1&2&1&1&0&0 \\ 0&1&2&-2&1&0 \\ 0&0&1&-3&1&1 \end{array}\right)
    \xrightarrow{\text{hacia arriba}}
    \left(\begin{array}{ccc|ccc} 1&0&0&-4&1&3 \\ 0&1&0&4&-1&-2 \\ 0&0&1&-3&1&1 \end{array}\right)
  \]
  \vfill
  \[
    A^{-1} = \begin{pmatrix} -4 & 1 & 3 \\ 4 & -1 & -2 \\ -3 & 1 & 1 \end{pmatrix}. \quad \text{Verificación: } A A^{-1} = I_3. \ \checkmark
  \]
\end{frame}

\begin{frame}{Ejercicio 8: Invertibilidad con Parámetro $k$}
  Sea $A = \begin{pmatrix} 1 & 1 & 1 \\ 1 & k & 2 \\ 2 & 3 & k \end{pmatrix}$. Reduciendo $[A \mid I_3]$:
  \[
    \left(\begin{array}{ccc|ccc} 1&1&1&1&0&0 \\ 0&1&k-2&-2&0&1 \\ 0&0&-k^2+3k-1&\ast&\ast&\ast \end{array}\right)
  \]
  \vfill
  \begin{itemize}
    \item $A$ es invertible $\iff -k^2+3k-1 \neq 0 \iff k \neq \frac{3 \pm \sqrt{5}}{2}$.
    \item Para $k=0$: $A^{-1} = \begin{pmatrix} -6 & 3 & 2 \\ 4 & -2 & -1 \\ 3 & -1 & -1 \end{pmatrix}$.
    \item Solución de $A\vv{x} = (1, 0, 4)^T$: $\vv{x} = A^{-1}\vv{b} = (2, 0, -1)^T$.
  \end{itemize}
\end{frame}

\begin{frame}{Ejercicio 9: Ecuación Matricial $A X B = C$}
  Resolver $AXB = C$ despejando $X = A^{-1} C B^{-1}$:
  \begin{itemize}
    \item $A = \begin{pmatrix} 2 & 1 \\ 1 & 1 \end{pmatrix} \implies A^{-1} = \begin{pmatrix} 1 & -1 \\ -1 & 2 \end{pmatrix}$.
    \item $B = \begin{pmatrix} 1 & 3 \\ 0 & 1 \end{pmatrix} \implies B^{-1} = \begin{pmatrix} 1 & -3 \\ 0 & 1 \end{pmatrix}$.
    \item $Y = A^{-1} C = \begin{pmatrix} 1 & -1 \\ -1 & 2 \end{pmatrix} \begin{pmatrix} 3 & 5 \\ 2 & 4 \end{pmatrix} = \begin{pmatrix} 1 & 1 \\ 1 & 3 \end{pmatrix}$.
    \item $X = Y B^{-1} = \begin{pmatrix} 1 & 1 \\ 1 & 3 \end{pmatrix} \begin{pmatrix} 1 & -3 \\ 0 & 1 \end{pmatrix} = \begin{pmatrix} 1 & -2 \\ 1 & 0 \end{pmatrix}$.
  \end{itemize}
  \vfill
  Comprobación: $A X B = \begin{pmatrix} 3 & 5 \\ 2 & 4 \end{pmatrix} = C$. \checkmark
\end{frame}

\begin{frame}{Ejercicio 10: Aplicación en Ing. Química (Tanques)}
  Balance de masa $M\vv{c} = \vv{s}$ con $M = \begin{pmatrix} 3 & -1 & 0 \\ -1 & 4 & -2 \\ 0 & -2 & 3 \end{pmatrix}$:
  \begin{enumerate}
    \item Matriz de respuesta por $[M \mid I_3] \to [I_3 \mid M^{-1}]$:
          \[
            M^{-1} = \frac{1}{21} \begin{pmatrix} 8 & 3 & 2 \\ 3 & 9 & 6 \\ 2 & 6 & 11 \end{pmatrix}.
          \]
    \item Para entrada $\vv{s}_1 = (10, 0, 5)^T$:
          $\vv{c}_1 = M^{-1}\vv{s}_1 = \begin{pmatrix} 30/7 \\ 20/7 \\ 25/7 \end{pmatrix} \approx \begin{pmatrix} 4.29 \\ 2.86 \\ 3.57 \end{pmatrix}$ g/L.
    \item Para entrada $\vv{s}_2 = (0, 15, 10)^T$:
          $\vv{c}_2 = M^{-1}\vv{s}_2 = \frac{1}{21} \begin{pmatrix} 65 \\ 195 \\ 200 \end{pmatrix} \approx \begin{pmatrix} 3.10 \\ 9.29 \\ 9.52 \end{pmatrix}$ g/L.
  \end{enumerate}
\end{frame}

\end{document}
